% sec-10: Document 10, parking and participant transportation standards.
\docpart{Parking and Transport}{PDAC LLM Oncology Clinical Trial Site Parking
and Patient Transportation Standards}

\section{Scope and Applicability}

\draftnote{State what this document governs and how it relates to the San Diego
Municipal Code parking requirements in document 05. Where the city minimum and
the trial requirement differ, the higher governs; say so once here.}

\section{Stall Count, Mix, and Dimensions}

\begin{mvtable}
\begin{tabularx}{\textwidth}{>{\raggedright\arraybackslash}p{3.4cm} >{\raggedright\arraybackslash}p{1.6cm} >{\raggedright\arraybackslash}p{2.6cm} >{\raggedright\arraybackslash}X}
\headrow \thead{Stall class} & \thead{Count} & \thead{Dimension} & \thead{Basis for the count}\\
\midrule
Participant & to be fixed & to be fixed & to be fixed \\
Accessible & to be fixed & to be fixed & to be fixed \\
\end{tabularx}
\end{mvtable}
\tabcapl{tab:parking}{The stall schedule. Stage 2 completes it and the count
column sums to the total fixed in the site parameter table in the front matter.}

\draftnote{Derive the participant stall count from the visit schedule in
\mvfile{trial-protocol/} rather than from a floor area ratio: a Phase 1 site with
up to eighteen treated participants has a peak concurrency that can be counted,
and counting it is more defensible than applying a general medical office ratio.}

\section{Accessible Parking and Path of Travel}

\draftnote{Accessible stall count and the accessible route, against the 2010 ADA
Standards for Accessible Design and Title 24. State the path of travel from the
accessible stall to the clinical zone entrance and keep it identical to the
wayfinding path in document 09 \sect{}6.}

\section{Drop-Off, Ambulance, and Loading Geometry}

\draftnote{Turning radius, canopy clearance, and the separation of the ambulance
approach from the participant drop-off. State the clearance that a transfer
ambulance needs and cite document 11 \sect{}2 for the evacuation route it
shares.}

\section{Post-Procedure Transportation}

\draftnote{A participant discharged after a pancreaticoduodenectomy does not
drive. State the requirement, the responsible role from document 14, the record
that proves the arrangement was made, and what happens when it fails. Take the
participant-facing framing from paper 18 of the deck table.}

\section{Transportation Demand Management and Monitoring}

\draftnote{The measures the city would expect and the monitoring report that
demonstrates them. Keep the reporting interval identical to document 05
\sect{}6.}
